\section{RAG 基本概念}

\subsection{RAG 整体框架}

RAG（Retrieval-Augmented Generation）是一个非常实用、非常基础的现代框架，无论工程还是科研都极为必要。

\begin{importantbox}{RAG 核心流程}
\begin{enumerate}
    \item 用户发出一个 \textbf{Query}（问题）
    \item Query 与外部数据的 \textbf{Index} 进行匹配——外部数据指非语言模型之外的、存在硬盘上的数据（结构化数据库、非结构化文档等）
    \item 对外部数据进行 \textbf{Chunking}（分块），每个 Chunk 创建一个 \textbf{Embedding Vector}
    \item 对 Query 也创建 Embedding Vector
    \item 计算 Query 与各 Chunk 之间的\textbf{相似度}（点积/余弦），取 Top-K 最相关的 Chunk
    \item 将相关 Chunk 作为 Context，与 Query 一起送入 LLM，生成回答
\end{enumerate}
\end{importantbox}

\subsection{两个核心 Model}

RAG 系统中需要两个模型：
\begin{enumerate}
    \item \textbf{Embedding Model}：将文本转换为向量表示（默认 text-embedding-ada-002，1536 维）
    \item \textbf{Language Model}：基于 Context 生成回答（默认 GPT-3.5 Turbo）
\end{enumerate}

\begin{knowledgebox}{RAG 的核心价值}
RAG 要求 LLM 仅基于检索出的 Context Information 回答问题，而非使用模型自身已有的知识——降低幻觉，提高对事实性要求高的场景的可靠性。
\end{knowledgebox}

\subsection{本章小结}

RAG 的本质是"检索增强生成"：先从外部数据中检索最相关的信息，再让 LLM 基于这些信息回答问题，从而减少幻觉、提高准确性。

